\scanpage{99}
\uline{Th} Soient $b,c,b',c',t\in k$, d'où $M=M(b,c,b',c',t)$, avec base can. $\underline{1},\underline{u},\underline{v},\underline{u}\,\underline{v}$, et formes $\operatorname{Tr}$ et $\det$ \ldots{} Pour que $M$ soit Azumaya, il faut et suffit que la forme $\det$ sur $\unsure{1k|k}\oplus k\underline{u}\oplus k\underline{v}$ soit non dégénérée, i.e. que son discriminant \unsure{arithmét|réduit} \unsure{et|ci} (précédemment désigné par $-Q$
\[
\Delta(b,c,b',c',t)=-(t^2-bb't+b^2c'+b'^2c-4cc')
\]
soit inv.

En termes plus intrinsèques, soit
\[
\deleted{M\ni}\qquad M\supset A_1,\deleted{\unsure{B}} A_2
\]
une algèbre \inserted{ass. unit. loc.} libre de rang $4$ sur $k$, \deleted{\unsure{sous une algèbre}} et deux sous-alg. $A_1,A_2$ \inserted{loc. libres} de rang $2$ \inserted{transverses}. On suppose que \deleted{\unsure{l'on peu}} dans $M$ il y ait des formes $\operatorname{Tr},\det$ avec les propr. hab, \inserted{et que $A_1,A_2$ ne \uuline{sont} pas (sur aucune fibre) contenues dans une sous-alg. de rg $3$ ($A_1+A_2$).} (ce qui implique que si on \unsure{décrit} sur $A_1$ gén. $u$, $A_2$ gén. $\unsure{v|V_2}$, d'où $b,\unsure{c|k},b',c',t$, on a $M\simeq M(b,c,b',c';t)$). Alors $M$ est azumaya sss la restriction de la forme \unsure{quad.|quadr.} $\det$ sur $A_1+A_2$ est non dégénérée,
% archive sheet number: 98
